\chapter{Inteligencia Artifical}

Considero que integrar un chat con IA al que preguntarle cosas del contexto
del libro puede ayudar mucho, estaria limitado poco mas del contexto, con un
Fine Tuning podria ser muy bueno. Sin embargo, hay un par de detalles que
debo solucionar antes y que no tenia contemplados

\section{ Preparativos }

Para que el bot pueda conocer el contexto, debe saber que tema estudias, debe
tener nocion de la posicion del libro, sin embargo, no puedo medir segun cuanto
scroll tiene su navegador porque si cambio de dispositivo, el scroll cambia
y pierdo la referencia. La solucion que se me ocurre es que debe saber que
capitulo, sección o sub sección estas leyendo y para eso, necesito referencias
y estas se logran con un identificador (id) que solo existe para los capitulos.

\subsection{ Contadores }

Ya que los capitulos tienen un contador unico, una variable numerica que voy
incrementando es suficiente, pero las subsecciones no tienen ese identificador
por que no van numeradas. Poner otro contador no es la solucion completa porque
incrementar el contador padre deberia reiniciar el contador hijo y este a su
vez debe reiniciar a sus hijos.

Como deben tener logica, deben ser objetos.

\subsection{ La posicion }

Ahora que ya hay contadores, ya identificamos elementos, entonces ahora podemos
guardar la posicion en la que nos encontramos, un capitulo, una sección, etc.
Aun no tengo un servicion de servidor activo porque incluir IA aun no lo
requiere pero puedo retomar posicion de lectura en el dispositivo actual, ya
es algo.

\subsection{ IA }

Ahora que ya se guarda la referencia de que estas leyendo, ahora toca integrar
la IA, tal vez se llame M(ath) I(nteligence) A(rtificial) - MIA. Empezaré a
entrenarla y establecer limites. Sin emabrgo, por mucho que esta función la
logre integrar, me temo que liberarla a la comunidad de momento estaría fuera
de mi alcancé pues mi servidor no cuenta con los recursos para tener una IA
ejecutandose, alrededor de 50 estudiantes simultaneos podrian saturar el
servicio y quizá no colapse pero se sentiria lento.